\documentclass[12pt]{article}
\usepackage[utf8]{inputenc}
\usepackage[T1]{fontenc}
\usepackage[margin=1in]{geometry}
\usepackage{setspace}
\usepackage{xcolor}
\usepackage[colorlinks=true, urlcolor=blue!60!black, linkcolor=blue!60!black]{hyperref}
\setlength{\parindent}{0em}
\setlength{\parskip}{0.6em}
\renewcommand{\baselinestretch}{1.0}
\begin{document}
\begin{center}
{\LARGE \textit{Myers v. United States}, 272 U.S. 52 (1926)\par}
\vspace{0.5em}
{\large Case Brief --- Two-Page Edition\par}
\end{center}

\section{Memory Jogger}
The President has exclusive power to remove executive officers he appoints with Senate consent; Congress may not require Senate approval of removals (\href{https://supreme.justia.com/cases/federal/us/272/52/}{272 U.S. 52}).

\section{Detailed Case Facts}
President Wilson appointed Frank Myers Portland's first-class postmaster for a four-year term under an 1876 statute providing that such postmasters may be removed by the President with the Senate's advice and consent. In 1920 Wilson had Myers removed without seeking Senate consent. Myers protested and sued for his lost salary.

\section{Procedural History}
The Court of Claims ruled for the government on laches; the Supreme Court rejected laches, heard reargument with Senator Pepper defending the statute, and affirmed on constitutional grounds.

\section{Judicial Votes}
6--3. Chief Justice Taft wrote for Justices Van Devanter, Sutherland, Butler, Sanford, and Stone. Justices Holmes, McReynolds, and Brandeis each dissented.

\section{Holding}
Article II gives the President exclusive power to remove executive officers appointed with Senate consent; the 1876 restriction, like the Tenure of Office Act of 1867, is unconstitutional.

\section{Analysis of Opinions}
Taft read Article II's Vesting and Take Care Clauses to require control over subordinates, treated the Senate's role in appointments as a narrow exception not extending to removals, and relied on the Decision of 1789 and decades of practice. Holmes said Congress, which creates the office, may condition removal. McReynolds found no express removal power and called tenure conditions legislative. Brandeis, surveying 137 years of practice, insisted that separation of powers exists to prevent arbitrary power, not to promote efficiency. \textit{Humphrey's Executor} later narrowed \textit{Myers}; \textit{Seila Law} and \textit{Trump v.\ Slaughter} (2026) revived it.

\section{Examples: Future Applications}
\textbf{Same-Side:} for-cause protection for commissioners wielding executive rulemaking and enforcement power is invalid after \textit{Slaughter}. \textbf{Opposite-Side:} a court-martial requirement for a naval officer appointed by the Secretary is valid under \textit{Perkins}. \textbf{Fence-Sitter:} single-layer protection for an inferior officer appointed by an agency head.

\section{Critique}
Historians argue that Taft overstated the consensus of 1789, and Justice Gorsuch has noted that plenary removal combined with broad delegation concentrates great power in one executive.

\section{Key Quotations}
\begin{quote}
``The doctrine of the separation of powers was adopted by the convention of 1787, not to promote efficiency but to preclude the exercise of arbitrary power. The purpose was not to avoid friction, but, by means of the inevitable friction incident to the distribution of the governmental powers among three departments, to save the people from autocracy.''
\end{quote}
\noindent --- Justice Brandeis, dissenting, 272 U.S.\ at 293.

\end{document}
